%% CV - Dalfin Rafael Solihin — Avelios Medical Working Student ML
%% Compile with: lualatex main_avelios_ml.tex

\documentclass[11pt,a4paper,sans]{moderncv}
\moderncvstyle{banking}
\moderncvcolor{blue}
\renewcommand*{\namefont}{\fontsize{34}{36}\bfseries\upshape}
\colorlet{firstnamecolor}{color1}
\colorlet{lastnamecolor}{color1}
\colorlet{namecolor}{color1}
\renewcommand*{\sectionstyle}[1]{{\sectionfont\color{color1}#1}}
\usepackage[utf8]{inputenc}
\AtEndPreamble{\hypersetup{colorlinks=true,linkcolor=blue,filecolor=magenta,urlcolor=blue,pdftitle={Dalfin Rafael Solihin - CV},pdfpagemode=UseNone,}}
\usepackage[scale=0.77]{geometry}
\usepackage{import}
\usepackage{needspace}

\name{Dalfin Rafael}{Solihin}
\address{Leipartstr. 28, 81369 Munich, Germany}{}{}
\phone[mobile]{+49 155 664 752 60}
\email{dalfinrafa@gmail.com}
\extrainfo{\href{https://www.linkedin.com/in/dalfin-rafael-solihin}{LinkedIn}, \href{https://github.com/discretediode}{GitHub}}

\begin{document}

\makecvtitle

% ============================================================
%     PROFILE STATEMENT
% ============================================================

\vspace{6pt}
\small{Mechatronics Engineering student at Munich University of Applied Sciences with hands-on experience in Python ML/AI engineering, LLM pipeline development, and deep learning model training. Trained and evaluated a custom computer vision model using PyTorch-based YOLOv8, and built a production-grade LLM automation pipeline integrating the Google Gemini API. Research assistantship at HM Munich reinforced systematic data analysis and Python scripting in a research environment. Excited to apply machine learning to real-world clinical workflows and contribute to the development of KIS 2.0 at Avelios Medical.}

% ============================================================
%     CORE COMPETENCIES
% ============================================================

\section{Core Competencies}
\vspace{1pt}
\begin{itemize}

\item \textbf{Python \& ML Engineering}: Python, NumPy, PyTorch, Ultralytics YOLOv8, scikit-learn, data pipeline design

\item \textbf{Computer Vision \& Deep Learning}: Custom model training, evaluation and benchmarking, image classification, YOLOv8/PyTorch

\item \textbf{LLM Integration \& Automation}: Google Gemini API, prompt engineering, LLM-driven data pipelines, PII handling

\item \textbf{Data Analysis \& Research}: Experimental measurement systems, data acquisition, structured reporting, Python scripting

\item \textbf{Technical Documentation}: Research-grade documentation, structured reporting, version control with Git/GitHub

\end{itemize}

% ============================================================
%     PROFESSIONAL EXPERIENCE
% ============================================================

\section{Professional Experience}
\vspace{3pt}
\begin{itemize}

\needspace{5\baselineskip}
\item{\cventry{Apr 2026 - Jul 2026}{Research Assistant}{Munich University of Applied Sciences}{Munich, Germany}{}{\vspace{1pt}
\begin{itemize}
    \item Conducted SFRA measurements on rotating electrical machines to detect short-circuits and earth faults, managing a multi-step experimental pipeline from setup to documented output
    \item Built electronic circuits and data acquisition setups for experimental testing using oscilloscopes, multimeters, and SFR analysers
    \item Developed a custom Python script to automate analysis and structured documentation of measurement data
    \item Collaborated with researchers to optimise testing procedures and improve experimental setups
\end{itemize}}}

\end{itemize}

% ============================================================
%     PROJECT EXPERIENCE
% ============================================================

\section{Project Experience}
\vspace{3pt}
\begin{itemize}

\needspace{5\baselineskip}
\item{\cventry{2025}{YOLOv8 Object Detection Model}{\href{https://github.com/discretediode}{github.com/discretediode}}{Python, PyTorch, Ultralytics}{}{\vspace{1pt}
\begin{itemize}
    \item Trained and evaluated a custom computer vision model using PyTorch-based YOLOv8 (Ultralytics) for image-based object detection and classification
    \item Configured training pipelines, tuned hyperparameters, and assessed model performance metrics across validation datasets
    \item Gained direct hands-on experience with the deep learning model development cycle: dataset preparation, training, evaluation, and iteration
\end{itemize}}}

\vspace{3pt}

\needspace{5\baselineskip}
\item{\cventry{2025}{AI Application Generator (CV Automator)}{\href{https://github.com/discretediode}{github.com/discretediode}}{Python, Google Gemini API, xhtml2pdf}{}{\vspace{1pt}
\begin{itemize}
    \item Built a local LLM automation tool integrating the Google Gemini API, engineered for production-grade document generation with PII stripping and re-injection
    \item Designed a multi-stage data pipeline covering input processing, LLM inference, and automated PDF assembly via xhtml2pdf
    \item Developed entirely using \textbf{Claude Code}, making AI-assisted development an everyday engineering practice
\end{itemize}}}

\vspace{3pt}

\needspace{5\baselineskip}
\item{\cventry{2025}{STM32 Embedded State Machine}{\href{https://github.com/discretediode}{github.com/discretediode}}{C, STM32, Embedded Systems}{}{\vspace{1pt}
\begin{itemize}
    \item Designed a timer-driven finite state machine in C on STM32 microcontroller with precise hardware debounce logic
    \item Demonstrated low-level systems programming and embedded engineering depth complementing software ML work
\end{itemize}}}

\end{itemize}

% ============================================================
%     EDUCATION
% ============================================================

\section{Education}
\vspace{1pt}
\begin{itemize}

\needspace{5\baselineskip}
\item{\cventry{Oct 2024 - present}{Bachelor in Mechatronics}{Munich University of Applied Sciences (HM Munich)}{Munich, Germany}{}{\vspace{1pt}
In progress, expected 2027/2028. Topics: embedded systems, control engineering, electronics, mechanical design, programming.
}}

\vspace{3pt}

\needspace{5\baselineskip}
\item{\cventry{Mar 2024 - Jul 2024}{Feststellungsprüfung, Technik (T/TI)}{Internationales Studienkolleg, Hochschule Kaiserslautern}{Kaiserslautern, Germany}{}{\vspace{1pt}
Overall grade: 1.5. Deutsch (2.0), Mathematik/Informatik (1.7), Physik (1.0). Equivalent to DSH-3/C1.
}}

\end{itemize}

% ============================================================
%     EXTRACURRICULAR
% ============================================================

\section{Extracurricular}
\vspace{1pt}
\begin{itemize}

\item{\cventry{Jan 2026 - present}{Active Member --- Aeronautics Division}{AkaModell München}{Munich, Germany}{}{\vspace{1pt}
\begin{itemize}
    \item Wing aerodynamic optimisation using MATLAB; structural design using Catia V5; hands-on electronics assembly and manual soldering
\end{itemize}}}

\end{itemize}

% ============================================================
%     LANGUAGES
% ============================================================

\section{Languages}
\vspace{1pt}
\begin{itemize}
\item {Indonesian (native), English (fluent), German (fluent --- DSH-3/C1), Chinese (advanced).}
\end{itemize}

% ============================================================
%     REFERENCES
% ============================================================

\section{References}
\vspace{1pt}
\begin{itemize}
\item{Available upon request.}
\end{itemize}

\end{document}
